\begin{abstract}
  Modern security relies heavily on public key cryptography,
  and it is necessary to protect private keys, and secrets in general, from attackers.
  Against a highly-capable adversary it is ideal to store secrets outside of main memory,
  which is easy on general purpose systems with the now widely-available Trusted Platform Module (TPM) 2.0.
  However, the lack of integration between the TPM and the OS makes protecting secrets with automated availability needs difficult.
  Using only existing mechanisms available on Linux: SELinux, Integrity Measurement Architecture (IMA), Extended Verification Module (EVM), and Linux Unified Key Setup (LUKS),
  we develop a strategy to authenticate specific OS entities to protect TPM-stored secrets,
  without limiting access to the TPM in general.

  \iflncs\keywords{First keyword \and Second keyword \and Another keyword.}\fi
\end{abstract}

\ifacm
\begin{CCSXML}
<ccs2012>
<concept>
<concept_id>10002978.10002979.10002980</concept_id>
<concept_desc>Security and privacy~Key management</concept_desc>
<concept_significance>500</concept_significance>
</concept>
<concept>
<concept_id>10002978.10002991.10002993</concept_id>
<concept_desc>Security and privacy~Access control</concept_desc>
<concept_significance>500</concept_significance>
</concept>
<concept>
<concept_id>10002978.10002991.10002992</concept_id>
<concept_desc>Security and privacy~Authentication</concept_desc>
<concept_significance>500</concept_significance>
</concept>
<concept>
<concept_id>10002978.10002991.10010839</concept_id>
<concept_desc>Security and privacy~Authorization</concept_desc>
<concept_significance>500</concept_significance>
</concept>
</ccs2012>
\end{CCSXML}

\ccsdesc[500]{Security and privacy~Key management}
\ccsdesc[500]{Security and privacy~Access control}
\ccsdesc[500]{Security and privacy~Authentication}
\ccsdesc[500]{Security and privacy~Authorization}

\keywords{Trusted Platform Module, Measured Boot, SELinux, Integrity Measurement Architecture} % TODO finalize
\fi

